% ===================================================================
\section{Biology, ecological thinking and the wildebeest (EE01)}
% ===================================================================

\subsection{Why biology resists a few neat equations}
Lecture 1 opens with three complaints: biology needs a lot of memorising, it is description-heavy, and it seems
to have no structure of ``a few equations from which the rest can be derived''. The answer the course gives is not
a defence of biology but a description of what biological systems are like. Slide 15 lists four features:

\begin{enumerate}
\item \textbf{Diversity and interconnectedness}, within a level (many species in one community) and across levels
(a change in a muscle protein can change a herd; a herd can change a river).
\item \textbf{Stochasticity}: which calf is caught, which seed lands where, which mutation arises.
\item \textbf{Historicity}: the present state depends on the particular path that produced it.
\item \textbf{Emergence}: properties of the whole are not found by dissecting the parts (``reductionism often fails'').
\end{enumerate}

\begin{chain}[: from the four features to the course's method]
Each feature blocks the physicist's route in its own way.
\begin{itemize}
\item Interconnectedness $\Rightarrow$ there is no isolated ``system'' to write equations for, and cutting a
system off from its connections changes how it behaves.
\item Stochasticity $\Rightarrow$ even a correct model predicts only probabilities.
\item Historicity $\Rightarrow$ the ``initial conditions'' are themselves a historical fact that has to be described,
not derived.
\item Emergence $\Rightarrow$ the behaviour of the whole cannot be assembled from the laws of the parts.
\end{itemize}
Therefore (slide 16): \textbf{mathematical formulation is typically difficult}, and when it can be done
\textbf{the relation with reality is statistical and conditional}. Two responses are possible: describe what is
there, or build simplified models with caveats. BI2113 uses (1) \emph{oversimplified mathematical models},
(2) \emph{a focus on concepts}, (3) \emph{historical narrative}, (4) \emph{examining the pedagogy}, and
accepts a trade-off: \textbf{breadth over depth}.

This also answers the three complaints: memorising and description come from diversity and history; the lack of
structure comes from mathematics working only locally, with caveats, while concepts and narrative carry the
connecting threads.
\end{chain}

\begin{keyidea}
``Oversimplified'' is a method, not an apology. A model is judged by what it teaches, not by how realistic it is:
the logistic equation is wrong for every real population and still explains why populations level off. The
course's wrap-up lecture repeats this: \emph{never mistake your model for reality}, and a model that shows the
expected behaviour is not thereby the right model.
\end{keyidea}

\subsection{Levels of organisation and the domain of the course}
Slide 13 draws a ladder: biomolecules $\to$ cells $\to$ tissues $\to$ organs $\to$ individual $\to$
population $\to$ community $\to$ ecosystem $\to$ biosphere. The levels from population upward change on
ecological time scales. \textbf{Changed trait values in a population are microevolution}; accumulated, they lead
through \textbf{speciation} to \textbf{macroevolution}. The individual's life cycle (meiosis to a gamete,
fertilisation to a cell) links the sub-organismal levels back to the individual. Slide 14 removes the
sub-organismal levels: \textbf{the domain of BI2113 is the individual and above}.

Different levels run on different clocks: muscle contraction in milliseconds, a migration in a year, the rise of a
lineage in millions of years. A model that follows one clock can treat the slower processes as fixed and the
faster ones as instantaneous. That separation of time scales is what makes simple models possible at all.

\subsection{Proximate and ultimate questions}
Slide 5 splits the wildebeest migration into \emph{how} and \emph{why}:
\begin{itemize}
\item \textbf{How} (proximate): the mechanism inside one individual now: physiology, development, cues, behaviour rules.
\item \textbf{Why} (ultimate): the benefit that made selection favour the trait; its evolutionary history.
\end{itemize}
Both kinds of answer apply to \emph{every} trait. They are not rival explanations; they answer different
questions. (Mayr's names for them, and his example of the warbler that migrates on day-length cues while the
screech owl in the same wood stays, are background reading.)

\begin{trap}
Confusing the \emph{mechanism} of synchronous calving with its \emph{benefit}. How females synchronise (a fall in
forage protein triggers the cycle; a silent first ovulation cues bulls and other females) is proximate. Why it pays
(swamping predators) is ultimate. ``Calves stand within 6 minutes'' is proximate; ``this lets them keep up with a
herd that must move'' is ultimate.
\end{trap}

\subsection{The wildebeest case}
\subsubsection*{How: three adaptations}
\begin{enumerate}
\item \textbf{Efficient muscle.} 62.6\,\% efficient, against 41.8\,\% (cow) and 34\,\% (mouse) (Curtin et al.\
2018). Efficiency $\eta$ is the fraction of chemical energy used that comes out as mechanical work; the rest is
heat. For work $W$ the fuel burnt is $W/\eta$ and the heat released is $W(1/\eta-1)$. A cow therefore needs
$0.626/0.418=1.50\times$ the fuel, but releases $1.39/0.60=2.3\times$ the heat. Heat is what is left over, so it
grows much faster than fuel as $\eta$ falls; in a hot, dry habitat, shedding heat costs water.
\item \textbf{Synchronous births.} About 250{,}000 calves in about three weeks from mid-February, roughly 80\,\% of the
year's calves: about 12{,}000 a day, 500 an hour. Predators with a limited capacity to chase, kill and eat cannot
take more than a fixed number per day, so concentrating births lowers the \emph{fraction} lost
(\textbf{predator satiation}).
\item \textbf{Precocial calves.} A calf stands in about 6 minutes and walks in 30; topi and hartebeest calves take
much longer. A herd that must keep moving cannot afford helpless young.
\end{enumerate}

\begin{chain}[: why synchrony protects calves (a bound)]
Let $N$ calves be born over $w$ days and let predators kill at most $K$ per day. Losses in the window are at most
$Kw$, so the fraction lost is at most $\min(1,\,Kw/N)$. Shortening $w$ lowers the cap in proportion, as long as
$K<N/w$. If the herd is small ($N/w<K$) there is no cap, and a pulse of births may even attract predators: synchrony
works only with large numbers.
\end{chain}

\subsubsection*{Why: opposing gradients}
Rainfall rises from the south-east short-grass plains ($\sim$400\,mm) to the north-west woodland ($\sim$1200\,mm),
but soil fertility and grass protein run the other way (the plains sit on nutrient-rich volcanic ash). The puzzle
the slide sets: if the woodland is wetter all year, why do herds go onto the dry plains at the start of the rains?

\begin{chain}[: why a migration, rather than staying put]
\begin{enumerate}
\item Rain controls \emph{how much} green grass exists; fertility controls \emph{how good} it is.
\item The two gradients are opposed, so no single place is best all year: the plains are best (protein-rich, short
grass) while the rains last; the woodland when the plains dry out.
\item An animal that compares sites over a large enough distance therefore keeps moving. In Holdo et al.'s
models (background), rules using intake combined with plant nitrogen fit the observed herd distribution far better
than rules using biomass alone, and a migration appears only when the tracking radius exceeds about 80--100\,km;
with only local comparisons the simulated herd gets trapped.
\end{enumerate}
So ``they follow the rain'' is incomplete: they follow \emph{quality}, and the migration is an emergent property of
many animals following a simple rule on a landscape with opposed gradients.
\end{chain}

\subsubsection*{What the migration does to the ecosystem}
\begin{itemize}
\item \textbf{Death}: annual mass drownings at the Mara River crossings (about 6{,}000 a year) are a recurring pulse of
food and nutrients for fish and scavengers; bones release phosphorus for years.
\item \textbf{Fire}: grazing removes fuel. After rinderpest was eradicated in the 1960s the herd surged, fire
declined and trees recovered (background: a virus-to-landscape trophic cascade).
\end{itemize}

\subsubsection*{Why so many? (Hopcraft, background but a model of reasoning)}
Wildebeest outnumber the next herbivore (zebra) about six-fold. The argument eliminates explanations by comparison:
against topi and hartebeest (biology helps but is not enough: wildebeest are not that dense elsewhere); against
resident wildebeest herds (same biology, 30--240$\times$ fewer: residents are \emph{habitat-limited}); against zebra
(also migratory but far fewer: probably \emph{predation-limited}). The general recipe for migratory
superabundance: \textbf{refuge areas, an ephemeral high-quality resource, and connectivity} between patches.

\begin{numbers}
$\sim$1{,}500\,km a year; herd $\approx$1.2--1.3 million; muscle 62.6\,\% (cow 41.8, mouse 34); 250{,}000 calves in
$\sim$3 weeks ($\approx$80\,\% of the year's); calf stands in $\sim$6 min; rain 400\,$\to$\,1{,}200\,mm SE$\to$NW;
tracking radius $\gtrsim$80--100\,km; $\sim$6{,}000 drownings a year.
\end{numbers}

\begin{statements}
\begin{enumerate}[label=\arabic*.]
\item ``Wildebeest migrate because the plains get more rain early in the season.''
\F The woodland is wetter all year. The pull of the plains is grass \emph{quality} (protein, on fertile soil).
\item ``Proximate and ultimate explanations of a trait compete; at most one can be right.''
\F They answer different questions (how vs why) and both apply to every trait.
\item ``Synchronous calving lowers the number of calves predators take per day.''
\F It lowers the \emph{fraction} taken: predators are saturated, so the number per day is capped while births
per day are huge.
\item ``Because the wildebeest muscle is 1.5$\times$ more fuel-efficient than the cow's, it also releases 1.5$\times$ less heat.''
\F Heat per unit work is $1/\eta-1$: 0.60 against 1.39, a factor of 2.3. The heat ratio is bigger than the fuel ratio.
\item ``If biological systems are stochastic, a model of them cannot be useful.''
\F Models then predict probabilities and trends; the relation with data is statistical and conditional, which is
still informative.
\item ``Removing 6{,}000 drownings a year would change the wildebeest population a lot.''
\F About 0.5\,\% of a food-limited herd; density dependence soon offsets it. The river ecosystem, which loses a
nutrient subsidy, changes far more (2025 end-sem Q11).
\end{enumerate}
\end{statements}

\begin{longans}[: ``Why is the wildebeest the dominant herbivore of the Serengeti?'']
(1) \emph{Biology}: grazing adaptations for short protein-rich grass; efficient muscle for long travel; synchronous,
precocial calving that swamps and outruns predators. (2) \emph{Environment}: opposed rainfall and fertility
gradients, so high-quality grass exists somewhere reachable all year if animals move over $\gtrsim$80\,km.
(3) \emph{Comparisons}: residents with the same biology are habitat-limited; zebra, also migratory, are
predation-limited. (4) \emph{General recipe}: refuge, ephemeral high-quality resource, connectivity. (5) Note the
two-way interaction: the herd changes fire regimes and river nutrients.
\end{longans}
